\section*{अध्याय \ref{ch:b3:lab-techniques-3} --- प्रयोगशाला तकनीकें III: शोध-अभ्यास}

\begin{solution}{exo:b3:lab-techniques-3:1}
$(0.50/1013)^3 = \num{1.2e-10}$। रिसाव, काँच पर अधिशोषित गैस और ग्रीस तथा सेप्टम से गैस-निर्गमन
इससे कहीं अधिक छोड़ते हैं; अतिरिक्त चक्रों की तुलना में शुष्क, गर्म काँच के उपकरण और अच्छे जोड़
अधिक महत्त्वपूर्ण हैं।
\end{solution}

\begin{solution}{exo:b3:lab-techniques-3:2}
$\qty{20.0}{mmol}/\qty{1.48}{mol/L} = \qty{13.5}{mL}$।
\end{solution}

\begin{solution}{exo:b3:lab-techniques-3:3}\emergencystretch=3em
$120 + 10 \times 1.00782503 + 55.93493633 - 0.00054858 = \qty{186.01264}{u}$। त्रुटि: $10^6 \times (186.0129 - 186.01264)/186.01264 = +1.4$ ppm,
5 ppm के भीतर भली-भाँति।
\end{solution}

\begin{solution}{exo:b3:lab-techniques-3:4}
प्राथमिक: शोध-पत्रिका लेख, पेटेंट, शोध-प्रबंध। द्वितीयक: समीक्षा,
स्थिरांकों की पुस्तिका। तृतीयक: पाठ्यपुस्तक।
\end{solution}

\begin{solution}{exo:b3:lab-techniques-3:5}
$M = \qty{194.0}{g/mol}$: C $96.0/194.0 = 49.48$\,\%, H $10.0/194.0 = 5.15$\,\%, N $56.0/194.0 = 28.87$\,\%। अंतर 0.17, 0.07 और 0.16
अंक: 0.4 के भीतर, विश्लेषण सूत्र का समर्थन करता है।
\end{solution}

\begin{solution}{exo:b3:lab-techniques-3:6}
$\qty{170.0}{mg}/\qty{212.0}{g/mol} = \qty{0.802}{mmol}$; $c = \qty{0.802}{mmol}/\qty{0.54}{mL} = \qty{1.5}{mol/L}$ (दो सार्थक अंक, जो आयतन के
पाठ्यांक से तय होते हैं)।
\end{solution}

\begin{solution}{exo:b3:lab-techniques-3:7}
माध्य 73.0 और 77.0\,\%; दोनों मानक विचलन $\sqrt{10/4} = 1.58$; $s_p = 1.58$। $t = 4.0/(1.58\sqrt{2/5}) = 4.0$, जो
2.31 से अधिक है: 95\,\% स्तर पर सुधार सार्थक है।
\end{solution}

\begin{solution}{exo:b3:lab-techniques-3:8}
$F = (0.80/0.30)^2 = 7.1 > 5.05$: पहली विधि सार्थक रूप से कम परिशुद्ध है।
\end{solution}

\begin{solution}{exo:b3:lab-techniques-3:9}
संदमक-युक्त छोटी बोतलें ख़रीदिए, प्रत्येक पर प्राप्ति और खोलने की तिथि लिखिए, बंद,
अँधेरे और ठंडे स्थान पर रखिए; किसी भी आसवन या वाष्पन से पहले और खोलने के बाद
नियत अंतरालों पर पेरॉक्साइडों के लिए परीक्षण कीजिए; तिथि पार कर चुकी, या परीक्षण में धनात्मक,
बोतलों को सांद्रित किए बिना ख़तरनाक-अपशिष्ट मार्ग से निपटाइए।
\end{solution}

\begin{solution}{exo:b3:lab-techniques-3:10}
संकट: मुक्त हाइड्रोजन (आग, दाब), जल के साथ NaH की अभिक्रियाशीलता,
और सोडियम हाइड्राइड के साथ DMF का एक ज्ञात ऊष्माक्षेपी अपघटन जो गर्म करने पर
अनियंत्रित हो सकता है; DMF प्रजनन के लिए भी विषैला है। नियंत्रण: विलायक (THF या
2-मेथिलटेट्राहाइड्रोफ़्यूरान) या क्षारक को प्रतिस्थापित कीजिए; यदि वही रखें, तो पहले छोटे पैमाने पर, निम्न
ताप पर, ताप पर नज़र रखते हुए NaH अंशों में मिलाते हुए, धूम-हुड में कवच के
पीछे, बबलर से निकास करते हुए, शीतलन-कुंड तैयार रखकर चलाइए; लिखित
प्रक्रिया, एक दूसरा व्यक्ति उपस्थित, और अंत में सुरक्षा उपकरण।
\end{solution}

\begin{solution}{exo:b3:lab-techniques-3:11}\emergencystretch=10em
आपेक्षिक योगदान: 0.50, 0.50, $0.010/20.0 = 0.05$, $0.010/15.0 = 0.067$ और 0.10\,\%; संयुक्त
$\sqrt{0.25 + 0.25 + 0.0025 + 0.0044 + 0.010} = 0.72$\,\%। दोनों समाकल प्रभावी हैं: तुला की तुलना में बेहतर संकेत-रव अनुपात और
सावधानीपूर्ण आधार-रेखा तथा कला संशोधन अधिक महत्त्वपूर्ण हैं।
\end{solution}

\begin{solution}{exo:b3:lab-techniques-3:12}
$Q = (89.0 - 82.0)/(89.0 - 80.6) = 0.83 > 0.71$: परीक्षण उसे चिह्नित करता है। उसे अस्वीकार करने से पहले नोटबुक में
किसी कारण की जाँच कीजिए (ब्यूरेट का ग़लत पाठ्यांक, एक वायु-बुलबुला, अनुमापक का भिन्न बैच); यदि कोई
मिले, तो उसे अस्वीकार कीजिए और कारण बताइए; यदि नहीं, तो उसे रखिए, परीक्षण का प्रतिवेदन दीजिए, या
मापन दोहराइए।
\end{solution}

\begin{solution}{pb:b3:lab-techniques-3:1}
\textbf{1.} $\ce{FeCl2 + 2NaC5H5 -> Fe(C5H5)2 + 2NaCl}$।

\textbf{2.} सोडियम साइक्लोपेन्टाडाइएनाइड ऑक्सीजन और जल से नष्ट हो जाता है, और आयरन(II)
क्लोराइड नम वायु में आयरन(III) में ऑक्सीकृत हो जाता है।

\textbf{3.} $\qty{5.00}{g}/\qty{126.8}{g/mol} = \qty{0.0394}{mol}$; $\times \qty{185.8}{g/mol} = \qty{7.33}{g}$; लब्धि $5.86/7.33 = 80$\,\%।

\textbf{4.} $(0.10/1000)^3 = 10^{-12}$ (सिद्धांततः)।

\textbf{5.} द्वितय को गर्म करके, जिसमें प्रतीप डील्स--ऐल्डर अभिक्रिया होती है, और
वाष्पशील एकलक को आसवित करके; कक्ष ताप पर एकलक डील्स--ऐल्डर अभिक्रिया द्वारा फिर से
द्वितयित हो जाता है, इसलिए उसे ठंडा रखा और तुरंत प्रयुक्त किया जाता है।

\textbf{6.} प्रत्येक अंश हाइड्रोजन और ऊष्मा मुक्त करता है: उसे धीरे-धीरे मिलाने से गैस-प्रवाह
और ताप नियंत्रण में रहते हैं।

\textbf{7.} \qty{186.0126}{u}; त्रुटि $+1.4$ ppm।

\textbf{8.} C $120.0/185.8 = 64.58$\,\%, H $10.0/185.8 = 5.38$\,\%; प्राप्त 64.41 और 5.47: अंतर $-0.17$ और $+0.09$, 0.4 के भीतर।

\textbf{9.} दसों हाइड्रोजन सममिति से तुल्य हैं, और वलय तेज़ी से घूमते हैं।

\textbf{10.} एक ही संकेत, क्योंकि दसों कार्बन तुल्य हैं।

\textbf{11.} एकल-क्रिस्टल X-किरण संरचना।

\textbf{12.} नहीं: 18-इलेक्ट्रॉन संकुल के रूप में (\cref{ch:b3:organometallic-bonding}) यह वायु में स्थायी है, अपने
पूर्वगामियों के विपरीत।

\textbf{13.} इसके संकेत फ़ेरोसीन के संकेतों पर अतिव्याप्त नहीं होते, यह एक शुद्ध, अवाष्पशील,
अनार्द्रताग्राही ठोस है जिसे सटीकता से तौला जा सकता है, और यह नमूने के प्रति
अक्रिय है।

\textbf{14.} स्कैनों के बीच प्रत्येक संकेत की पूर्ण विश्रांति (कई $T_1$ का विलंब),
एकसमान उत्तेजन, पर्याप्त संकेत-रव अनुपात, और सावधानीपूर्ण कला तथा आधार-रेखा
संशोधन।

\textbf{15.} फ़ेरोसीन $\ce{C10H10Fe}$ \qty{185.8}{g/mol}; 1,3,5-ट्राइमेथॉक्सीबेंज़ीन $\ce{C9H12O3}$ \qty{168.0}{g/mol}।

\textbf{16.} $P_x = 3.942 \times (3/10) \times (185.8/168.0) \times (15.0/20.0) \times 0.999 = 0.980$: 98.0\,\%।

\textbf{17.} आपेक्षिक $\sqrt{0.50^2 + 0.50^2 + 0.05^2 + 0.067^2 + 0.10^2} = 0.72$\,\%; निरपेक्ष $0.72\,\% \times 98.0 = 0.7$ प्रतिशत-अंक।

\textbf{18.} हाँ, पर दुर्बल रूप से: समान संघटन की 2\,\% अशुद्धि C और H को उससे कम बदलती है
जितना विश्लेषण पकड़ सकता है; qNMR अधिक सूचनाप्रद शुद्धता-मापन है।

\emergencystretch=3em \textbf{19.} सोडियम हाइड्राइड (जल के साथ हाइड्रोजन, आग), निकलने वाली हाइड्रोजन, THF और
साइक्लोपेन्टाडाईन (अत्यधिक ज्वलनशील; THF \omterm{def:b3:lab-techniques-3:peroxide}{पेरॉक्साइड-निर्माता}), द्वितय का गर्म
विभंजन, आयरन(II) क्लोराइड (उत्तेजक), और फ़ेरोसीन (ज्वलनशील ठोस, हानिकारक)।

\textbf{20.} सोडियम हाइड्राइड: अक्रिय गैस के अधीन तेल-परिक्षेपण के रूप में तौला गया, विलोडित,
ठंडे किए गए विलयन में अंशों में मिलाया गया, गैस का निकास बबलर से, पास में कोई
जल नहीं। THF: किसी स्तंभ से शुष्क और पेरॉक्साइड-मुक्त, धूम-हुड में नाइट्रोजन के अधीन प्रयुक्त।

\textbf{21.} अक्रिय गैस के अधीन और शीतलन के साथ, धीरे-धीरे आइसोप्रोपेनॉल, फिर
एथेनॉल, फिर जल मिलाकर, हर बार गैस-निष्कासन रुकने तक प्रतीक्षा करते हुए।

\textbf{22.} एक उत्क्रमणीय एक-इलेक्ट्रॉन तरंग, $\ce{Fe(C5H5)2+}/\ce{Fe(C5H5)2}$, जिसके शिखर \qty{25}{\celsius} पर लगभग 59 mV दूर होते हैं
(\cref{ch:b3:electrode-kinetics}); यह युग्म तीव्र, स्थायी और विलायक के प्रति लगभग असंवेदी है,
जो इसे एक सुविधाजनक आंतरिक संदर्भ बनाता है।

\textbf{23.} नोटबुक: दिनांक, योजना, मात्राएँ और स्रोत, \omterm{def:b3:lab-techniques-3:assessment}{जोखिम मूल्यांकन}, जो
किया गया और देखा गया, लब्धियाँ और \omterm{def:b3:lab-techniques-3:notebook}{अपरिष्कृत आँकड़ों} की फ़ाइलों के नाम। \omterm{def:b3:lab-techniques-3:literature}{सहायक
सूचना}: पूर्ण प्रक्रिया और सभी अभिलक्षणन आँकड़े, स्पेक्ट्रमों की
प्रतियों सहित।

\textbf{24.} ताकि दूसरे, और स्वयं लेखक वर्षों बाद, उन्हें फिर से खोलकर पुनः प्रसंस्कृत कर सकें,
उनके पास जो भी सॉफ़्टवेयर हो: खोजने योग्य, सुलभ, अंतर्प्रचालनीय और पुनःप्रयोज्य।

\textbf{25.} फ़ेरोसीन नमूने की qNMR शुद्धता $98.0 \pm 0.7$\,\% (मानक अनिश्चितता) है।
\end{solution}
